\begin{table}[htbp]
\centering
\caption{Standardized Effect Sizes}\label{tab:sde}
\begin{threeparttable}
\small
\begin{tabular}{lcccccc}
\toprule
Outcome & $\hat{\beta}$ & SE & SD($Y$) & SDE & SE(SDE) & Classification \\
\midrule
\multicolumn{7}{l}{\textit{Panel A: Pooled}} \\
Log Price Index & -0.198 & 0.053 & 0.224 & -0.883 & 0.237 & Large negative \\
Log Rental Index & -0.068 & 0.019 & 0.246 & -0.278 & 0.078 & Large negative \\
Log Transactions & -0.563 & 0.232 & 0.632 & -0.890 & 0.367 & Large negative \\
\midrule
\multicolumn{7}{l}{\textit{Panel B: Heterogeneous (by ABSD intensity)}} \\
Log Price (Rounds 1--2, 10--15\%) & -0.060 & 0.033 & 0.224 & -0.270 & 0.148 & Large negative \\
Log Price (Round 5, 60\%) & -0.388 & 0.061 & 0.224 & -1.736 & 0.274 & Large negative \\
\bottomrule
\end{tabular}
\begin{tablenotes}[flushleft]\small
\item \textit{Notes:} \textbf{Country:} Singapore. \textbf{Research question:} Does escalating foreign-buyer stamp duty (ABSD, 0--60\%) suppress private residential property prices, transaction volumes, and rents in the Core Central Region relative to less foreign-exposed segments? \textbf{Policy mechanism:} Additional Buyer's Stamp Duty on foreign purchasers of private residential property, raised in five rounds from 10\% (December 2011) to 60\% (April 2023), increasing the upfront transaction cost and widening the wedge between foreign and domestic buyer reservation prices. \textbf{Outcome definition:} Log URA Property Price Index (hedonic, base 2009-Q1 = 100), log URA Rental Index, and log quarterly transaction count by market segment. \textbf{Treatment:} Binary (CCR high foreign-buyer exposure $\sim$16\% vs OCR/RCR low exposure $\sim$3--8\%). \textbf{Data:} URA via data.gov.sg, quarterly segment-level indices 2004-Q1 to 2025-Q4, 3 segments $\times$ 88 quarters = 264 observations. \textbf{Method:} Two-way FE DiD (segment + quarter), Driscoll-Kraay SEs (bandwidth = 4). \textbf{Sample:} Three non-landed private residential market segments (CCR, RCR, OCR); quarterly 2004--2025. SDE $= \hat{\beta} / \text{SD}(Y)$ where SD($Y$) is the pre-treatment (pre-2011-Q4) standard deviation of the log outcome for CCR. Classification refers to magnitude, not statistical significance: Large ($|$SDE$| > 0.15$), Moderate ($0.05$--$0.15$), Small ($0.005$--$0.05$), Null ($< 0.005$).
\end{tablenotes}
\end{threeparttable}
\end{table}

